\section{Technology: Connected but Dependent}
\label{sec:tech}

Almost all international data travel through submarine fibre-optic cables.
For a large coastal country the cable network is a mesh: dozens of
independent systems landing at many points, so a single fault is an
inconvenience. For a small island it is often a single thread. Tonga's
\tongaOutageDays{} days offline in 2022 were not caused by poor management
but by the arithmetic of redundancy.

\begin{figure}[t]
\centering
\begin{tikzpicture}
\begin{axis}[edge, height=0.6\linewidth,
    xmode=log, xlabel={Population, millions (log scale)},
    xmin=0.005, xmax=2000,
    xtick={0.01,0.1,1,10,100,1000}, xticklabels={0.01,0.1,1,10,100,1000},
    ymode=log, ymin=0.8, ymax=80,
    ytick={1,2,3,6,11,21,51}, yticklabels={0,1,2,5,10,20,50},
    ylabel={International submarine cables in service},
    legend pos=north west]
  \addplot[dotsBase, only if={other}{1}] table[x=pop_m,y=cables_p1] {\datadir/plot_countries.csv};
  \addlegendentry{Other coastal economies}
  \addplot[dotsSIDS, only if={sids}{1}] table[x=pop_m,y=cables_p1] {\datadir/plot_countries.csv};
  \addlegendentry{Small island states}
  \addplot[labelled=lab_cables] table[x=pop_m,y=cables_p1] {\datadir/plot_countries.csv};
\end{axis}
\end{tikzpicture}
\caption{\takeaway{Redundancy scales with size, not distance.} Cable counts
  rise with population (correlation of logs \corrCablesPop) but barely vary with
  remoteness (\corrCablesRemote). \nSIDSOneCable{} of \nSIDScables{} small
  island states, including \nPacificOneCable{} in the Pacific, have one
  international cable or none: a single cut takes them offline.
  \source{TeleGeography Submarine Cable Map (cables in service by 2025, excluding
  domestic systems); World Bank WDI.}}
\label{fig:cables}
\end{figure}


\Cref{fig:cables} counts international cables in service in each coastal
economy, using TeleGeography's public map data \parencite{telegeography2026}.
The striking pattern is that redundancy tracks \emph{size}: cable counts rise
steeply with population (correlation \corrCablesPop) and are only weakly
related to remoteness (\corrCablesRemote). Singapore has \cablesSingapore{}
international cables; the median island state \cablesSIDSMedian{}; the median
other coastal economy \cablesOtherMedian. \nSIDSOneCable{} of the
\nSIDScables{} small island states have at most one international cable. Cables
are lumpy, fixed-cost investments; small markets struggle to justify the
second one, which is precisely the one that matters when the first is cut.

Tonga now illustrates the remedy as well as the problem. In May 2026 a branch
of the trans-Pacific Hawaiki cable began service to Vava\okina u, financed with
Australian and New Zealand support, giving the kingdom a second international
link \parencite{aiffp2026tonga}.\footnote{TeleGeography records the branch
under the Hawaiki trunk's \enquote{ready for service} year, 2018; the Tongan
branch itself entered service in 2026. Tonga's count of \cablesTonga{} in
\cref{fig:cables} therefore reflects the 2026 position.}

\paragraph{Connected but dependent.} Connectivity has also changed what
dependence means. A small country's payments, government records, health
systems and schools increasingly run on services hosted abroad, priced in
foreign currency and governed by foreign law. That is mostly a gain---a tiny
administration can rent world-class software---but it means the recurring tax
now includes licence fees, bandwidth charges and the risk that a cut, a
sanction or a platform's change of terms switches services off.

\begin{counterpoint}[Leapfrogging]
Satellite constellations and mobile money cut directly against the thesis.
Low-Earth-orbit satellite internet reaches places where cables never will,
and mobile money has transformed financial access where banks were scarce:
in Kenya it lifted households' ability to share risk
\parencite{jack2014} and reduced poverty \parencite{suri2016}. Fast internet
measurably raised employment where it arrived in Africa
\parencite{hjort2019}.

The test is whether technology helps the edge \emph{faster than the centre}.
On access, it clearly does (\cref{fig:internet}). In 2000 the median
high-income country had \netRatioFirst{} times the share of people online of the
median island state; by \windowEnd{} the ratio was \netRatioLast. The
absolute gap peaked at \netGapPeak{} percentage points in \netGapPeakYear{} and
has since halved to \netGapLast. On \emph{intensity}, convergence is far
slower: fixed broadband lines per 100 people still differ by \bbGapLast{}
(\bbSidsLast{} against \bbHicLast), and in the Pacific the median is
\bbPacificLast. Leapfrogging is real for arrival; the gap persists in depth
of use, price and resilience. Satellite links also add a new dependency, on a
handful of foreign operators.
\end{counterpoint}

\begin{figure}[t]
\centering
\begin{tikzpicture}
\begin{axis}[edge, height=0.52\linewidth,
    xmin=1998, xmax=2023, xtick={2000,2005,2010,2015,2020},
    x tick label style={/pgf/number format/1000 sep={}},
    ymin=0, ymax=100, ylabel={Per 100 people},
    legend style={at={(0.02,0.98)},anchor=north west}]
  \addplot[lineSlate, mark=none] table[x=year,y=net_hic] {\datadir/internet_groups.csv};
  \addlegendentry{Internet users, high income}
  \addplot[lineTeal] table[x=year,y=net_sids] {\datadir/internet_groups.csv};
  \addlegendentry{Internet users, small island states}
  \addplot[lineSlate, dashed, dash pattern=on 4pt off 2pt, line width=1.1pt] table[x=year,y=bb_hic] {\datadir/internet_groups.csv};
  \addlegendentry{Fixed broadband, high income}
  \addplot[lineEmber] table[x=year,y=bb_sids] {\datadir/internet_groups.csv};
  \addlegendentry{Fixed broadband, small island states}
\end{axis}
\end{tikzpicture}
\caption{\takeaway{Access has converged; intensity has not.} In 2000 the
  median rich country had \netRatioFirst{} times the island share of people
  online; by \windowEnd{} it was \netHicLast\% against \netSidsLast\%. Fixed broadband, a proxy for heavy, reliable use, still trails
  by \bbGapLast{} lines per 100 people.
  Medians across countries reporting each year.
  \source{World Bank WDI (ITU data); UN M49 SIDS list.}}
\label{fig:internet}
\end{figure}

